% 第9章第6节：本章小结。

\section{本章小结}
\label{sec:opt-fn-summary}

章首问：连接$(0,0)$与$(1,1)$、使$\frac12\int_0^1u'^2dx$最小的曲线是哪一条，这个问题该怎样提、怎样解。现在可以按步骤回答。先明确候选空间与目标：取$X=H^1(0,1)$，能量在整个空间上有限，端点条件写进可行集$C$。再用微分计算小扰动的响应：方向导数是沿一条直线的一元导数，Gâteaux微分要求这些响应组成有界线性泛函，Fréchet微分进一步要求余项对所有小扰动统一地小；Hilbert空间中，微分借内积写成梯度，而梯度随内积改变。

最优条件按允许的扰动方向写出。无约束或仿射约束下微分在方向空间上为零，凸约束下变为变分不等式；这些都是必要条件，充分性来自凸性，或来自具有一致下界的二阶项。积分目标经分部积分分成内部与边界两部分：测试函数给出Euler--Lagrange方程或其弱形式，自由边界给出自然边界条件，固定边界则作为本质条件写进可行集。对章首问题，方程给出候选$u=x$，平方差展开证明它是唯一最小点。

这些条件都不负责存在性。最小化序列可以在函数值趋于下确界的同时振荡或逃逸，极限取不到最小值；直接法用有界性、弱收敛子列、弱序列闭性和弱下半连续性逐步排除这些情形，闭凸集投影则只用完备性与平行四边形恒等式给出存在唯一。最后，有限维近似或者改变候选集合（Ritz方法），或者改变目标（差分与图能量），二者都需要独立检查误差；固定函数的能量一致性不等于离散最小点收敛。表~\ref{tab:opt-fn-summary}汇总了本章各结论的条件与边界。具体的函数拟合由上一章主讲，时间演化、梯度流与方程解理论由第~\ref{part:dynamical-systems-control-and-decision}部分接续。

\begin{table}[!htb]
  \centering
  \small
  \begin{tabularx}{\linewidth}{@{}l>{\raggedright\arraybackslash}X>{\raggedright\arraybackslash}X@{}}
    \toprule
    问题 & 本章的工具与结论 & 所需条件与边界 \\
    \midrule
    小扰动的响应 & 方向导数、Gâteaux微分、Fréchet微分；Hilbert梯度 & 目标在点的开邻域内有限；Fréchet余项依赖所选范数；梯度依赖内积 \\
    候选是否最优 & 一阶条件、变分不等式、二阶条件 & 以最小点存在为前提；充分性需凸性或一致正定的二阶项 \\
    积分目标的条件 & Euler--Lagrange方程（弱形式）、自然边界条件 & 交换极限与积分需支配条件；点态方程需额外正则性 \\
    最小值是否达到 & 直接法；闭凸集投影存在唯一 & 下水平集有界、弱序列闭、弱下半连续；唯一性另需严格凸 \\
    怎样计算 & Ritz方法、差分能量、图能量 & 改变候选集合或目标；能量一致性不等于最小点收敛 \\
    \bottomrule
  \end{tabularx}
  \caption{泛函优化各环节的工具与条件。}
  \label{tab:opt-fn-summary}
\end{table}
